%%%%%%%%%%%%%%%%%%%%%%%%%%%%%%%%%%%%%%%%%%%%%%%%%%%%%%%%%%%%%%%%%%%%%%
% Lớp 10 | Hệ thức lượng trong tam giác | Giá trị lượng giác của góc (0 −180∘)
% Dạng 0H4?1-1: Xét dấu của biểu thức lượng giác
% Số câu: 27  (sinh tự động bởi tach_cau.py; từ nay thầy sửa trực tiếp file này)
%%%%%%%%%%%%%%%%%%%%%%%%%%%%%%%%%%%%%%%%%%%%%%%%%%%%%%%%%%%%%%%%%%%%%%

%%%%%%%%%%%%%%%% TRẮC NGHIỆM 4 PHƯƠNG ÁN %%%%%%%%%%%%%%%%

% nguồn: MaHoaNH_TracNghiem_Lop10
\begin{ex}%[0H4N1-1]%[ma:ddc31bf4]
	 Chọn khẳng định đúng trong các khẳng định sau
	\choice { $\sin(180^\circ-x)=-\sin x$}{ $\cos(180^\circ-x)=-\sin x$}{  $\tan(180^{\circ}-x)=\cot x$}{\True $\cos(180^{\circ}-x)=-\cos x$}
	\loigiai{ 
		Công thức  lượng giác cơ bản là $\cos(180^{\circ}-x)=-\cos x$.
	} 
\end{ex}

% nguồn: MaHoaNH_TracNghiem_Lop10
\begin{ex}%[0H4N1-1]%[ma:2110bea7]
	Cho góc $\alpha$ nhọn. Điều  khẳng định nào sau đây là \textbf{sai}?
	 \choice
		 {$\cos \alpha>0$}
		{$\tan \alpha>0$} 
		{\True  $\sin \alpha<0$}
		{$\cot  \alpha>0$}
	\loigiai{Góc nhọn có điểm biểu diễn thuộc góc phần tư thứ I, có giá trị $\sin \alpha >0$, $\cos \alpha >0$, $\tan  \alpha >0$ và $\cot \alpha>0$.}
\end{ex}

% nguồn: MaHoaNH_TracNghiem_Lop10
\begin{ex}%[0H4H1-1]%[ma:dd2eecb0]
	Trong các đẳng thức sau đây, đẳng thức nào \textbf{sai}? 
	 \choice
		{$\sin  0^\circ+\cos 0^\circ=1$}
		{$\sin 90^\circ+\cos 90^\circ=1$} 
		{$\sin  180^\circ+\cos 180^\circ=-1$}
		 {\True $\sin 60^\circ+\cos 60^\circ=1$}
	 \loigiai{
		 Ta có
		 \[\sin 60^\circ+\cos 60^\circ=\dfrac{\sqrt{3}}{2}+\dfrac{1}{2}\]
	} 
\end{ex}

% nguồn: MaHoaNH_TracNghiem_Lop10
\begin{ex}%[0H4Y1-1]%[ma:774aa9a8]
	Trong mặt phẳng $O x y$, véc-tơ pháp tuyến của đường thẳng $\Delta\colon x+2  y+3=0$ có toạ độ là
	 \choice
		{$(2;  3)$}
		{$(2;  1)$}
		 {\True$(1; 2)$}
		 {$(2;-1)$}
	 \loigiai{Véc-tơ pháp tuyến của đường thẳng $\Delta\colon x+2 y+3=0$ có toạ độ là $(1; 2)$.}
\end{ex}

% nguồn: MaHoaNH_TracNghiem_Lop10
\begin{ex}%[0H4Y1-1]%[ma:83dbf7f0]
	Đường  thẳng $\Delta\colon\heva{&x=1-t\\&y=3+2t}$. Véctơ chỉ phương có toạ độ là$\colon$
	 \choice
		{$\left(-1;3\right)$} 
		{\True  $\left(-1;2\right)$}
		{$\left(-2;1\right)$} 
		 {$\left(1;3\right)$}
	\loigiai{ Đường thẳng $\Delta$  có phương trình $\colon\heva{&x=1-t\\&y=3+2t}$ nên véctơ chỉ phương có toạ độ là$\colon\vec{u}\left(-1;2\right)$
	 }
\end{ex}

% nguồn: MaHoaNH_TracNghiem_Lop10
\begin{ex}%[0H4Y1-1]%[ma:823d2ee4]
	Trong mặt phẳng hệ toạ độ $Oxy$, một véc-tơ chỉ phương của đường thẳng $\heva{&x=3+5t \\ &y=-2+t}$ có tọa độ là 
	\choice 
		{$(3; -2)$} 
		 {\True $(5; 1)$}
		{$(-5;  1)$}
		{$(3;  2)$}
	\loigiai{ 
		Một véc-tơ chỉ phương của đường thẳng đã cho có tọa độ là  $(5; 1)$.
	} 
\end{ex}

% nguồn: MaHoaNH_TracNghiem_Lop10
\begin{ex}%[0H4B1-1]%[ma:f4359cd4]
	Trong mặt phẳng $O x y$, đường  thẳng $d\colon\heva{&x=1-2 t \\&y=-3+t}$, $(t \in \mathbb{R})$ đi qua điểm nào dưới đây?
	\choice 
		 {$B(1;-2)$}
		{$A(-2;  1)$}
		{\True $C(1;-3)$} 
		{$D(3;  1)$}
	 \loigiai{
		Với $t=0$,  ta có $\heva{&x=1-2\cdot 0 =1\\&y=-3+0=-3}$. Suy ra đường thẳng $d$ đi qua điểm $C(1;-3)$.
	} 
\end{ex}

% nguồn: MaHoaNH_TracNghiem_Lop10
\begin{ex}%[0H4Y1-1]%[ma:25b13782]
	Cho đường $(d)\colon \heva{&x=-1+2t\\&y=3-4t}  (t\in\mathbb{R})$. Véc-tơ nào sau đây là một véc-tơ chỉ phương của $(d)$?
	\choice 
		 {$\overrightarrow{a}=(1;2)$}
		{$\overrightarrow{a}=(-1;3)$} 
		{\True  $\overrightarrow{a}=(2;-4)$}
		 {$\overrightarrow{a}=(-1;-2)$}
	\loigiai{ 
		Véc-tơ chỉ phương của $(d)$ là  $\overrightarrow{a}=(2;-4)$.
	 }
\end{ex}

% nguồn: MaHoaNH_TracNghiem_Lop10
\begin{ex}%[0H4B1-1]%[ma:33a75a2a]
	Trong mặt phẳng tọa độ  $Oxy$, cho đường thẳng $d\colon 2x-3y+1=0$. Véc-tơ nào dưới đây là một véc-tơ pháp tuyến của đường thẳng $d$?
	 \choice
		{\True  $\overrightarrow{n}_1=(2;-3)$}
		{$\overrightarrow{n}_4=(3;2)$} 
		{$\overrightarrow{n}_2=(2;3)$} 
		 {$\overrightarrow{n}_3=(3;-2)$}
	 \loigiai{
		Đường thẳng $d\colon 2x-3y+1=0$ nhận $\overrightarrow{n}_1=(2;-3)$  làm một véc-tơ pháp tuyến.
	 }
\end{ex}

% nguồn: MaHoaNH_TracNghiem_Lop10
\begin{ex}%[0H4B1-1]%[ma:025a6998]
	Cho hàm  số $y=2 x^2-x-3$, điểm nào sau đây thuộc đồ thị hàm số?
	 \choice
		 {$M(1; 0)$}
		 {$M(2; 7)$}
		 {\True $M(0;-3)$}
		 {$M(-1;-2)$}
	 \loigiai{
		Với $x=0$, ta có $y=2\cdot0^2-0-3=-3$. Suy ra đồ thị hàm số đi qua điểm $M(0;-3)$. 
	} 
\end{ex}

% nguồn: MaHoaNH_TracNghiem_Lop10
\begin{ex}%[0H4Y1-1]%[ma:dc4ecbb9]
	Cho đường thẳng $(d)\colon 3x+2y-10=0$. Véc-tơ nào sau đây  là véc-tơ chỉ phương của $(d)$?
	 \choice
		 {$\vec{u}=(3;2)$}
		 {$\vec{u}=(-2;-3)$}
		{\True $\vec{u}=(2;-3)$} 
		{$\vec{u}=(3;-2)$} 
	 \loigiai{
		Đường thẳng $(d)\colon 3x+2y-10=0$  nhận $\vec{n}=(3;2)$ làm một véc-tơ pháp tuyến nên nhận $\vec{u}=(2;-3)$ làm một véc-tơ chỉ phương.
	} 
\end{ex}

% nguồn: MaHoaNH_TracNghiem_Lop10
\begin{ex}%[0H4Y1-1]%[ma:a7793b16]
	Cho đường thẳng $(d)$ có phương  trình $\heva{&x=1-t\\&y=3+2t}$. Khi đó, một véc-tơ pháp tuyến cùa đường thẳng $(d)$ là
	\choice 
		 {$\overrightarrow{n}=(-1;2)$}
		{$\overrightarrow{n}=(1;2)$} 
		{\True  $\overrightarrow{n}=(2;1)$}
		{$\overrightarrow{n}=(2;-1)$} 
	 \loigiai{
		Một véc-tơ pháp tuyến  cùa đường thẳng $(d)$ là $\overrightarrow{n}=(2;1)$.
	 }
\end{ex}

% nguồn: MaHoaNH_TracNghiem_Lop10
\begin{ex}%[0H4Y1-1]%[ma:2e9834f2]
	Cho đường thẳng $\Delta\colon x-3y-2=0$. Véctơ nào không  phải là véctơ pháp tuyến của $\Delta$.
	\choice 
		 {$(1;-3)$}
		{\True $(3;1)$} 
		 {$\left(\dfrac{1}{3};-1\right)$}
		 {$(-2;6)$}
	 \loigiai{
		Véctơ không phải là  véctơ pháp tuyến của $\Delta$ là $\vv{n} = (3;1)$.
	} 
\end{ex}

% nguồn: MaHoaNH_TracNghiem_Lop10
\begin{ex}%[0H4Y1-1]%[ma:a29069e2]
	Trong mặt phẳng  toạ độ $Oxy$, cho đường thẳng $d\colon 4x+3y-28=0$. Điểm nào sau đây thuộc đường thẳng $d$?
	\choice 
		 {$A(4; 0)$}
		{\True $C(7; 0)$} 
		{$B(-7;  0)$}
		{$D(3; 0)$} 
	 \loigiai{
		Vì $4\cdot 7+3\cdot 0-28=0$ nên $C(7; 0)\in  d$.
	} 
\end{ex}

% nguồn: MaHoaNH_TracNghiem_Lop10
\begin{ex}%[0H4Y1-1]%[ma:c4dbacb1]
	Một  véc-tơ pháp tuyến của đường thẳng $(d)\colon x-2 y+2023=0$ là
	 \choice
		{$\overrightarrow{n}=(1;2)$} 
		{$\overrightarrow{n}=(2;1)$} 
		 {$\overrightarrow{n}=(-2;1)$}
		 {\True $\overrightarrow{n}=(1;-2)$}
	\loigiai{ 
		 Một véc-tơ pháp tuyến của đường thẳng $(d)\colon x-2y+2023=0$ là $\overrightarrow{n}=(1;-2)$.
	 }
\end{ex}

% nguồn: MaHoaNH_TracNghiem_Lop10
\begin{ex}%[0H4Y1-1]%[ma:fc6d01b6]
	Trong  mặt phẳng toạ độ $Oxy$, cho đường thẳng $\Delta\colon 3x+4y+1=0$. Một véc-tơ pháp tuyến của đường thẳng $\Delta$ có tọa độ là
	\choice 
		{\True  $(3; 4)$}
		{$(3;  -4)$}
		 {$(4; -3)$}
		{$(-3;  4)$}
	 \loigiai{
		Một véc-tơ pháp tuyến của đường thẳng  $\Delta$ có tọa độ là $(3; 4)$.
	} 
\end{ex}

% nguồn: MaHoaNH_TracNghiem_Lop10
\begin{ex}%[0H4B1-1]%[ma:5e7725d3]
	Véc-tơ pháp tuyến của đường thẳng $x+2y-3=0$ là 
	\choice 
		{$\overrightarrow{n}=(1;-3)$} 
		{$\overrightarrow{n}=(2;-3)$} 
		 {$\overrightarrow{n}=(2;-1)$}
		{\True  $\overrightarrow{n}=(1;2)$}
	 \loigiai{
		Đường  thẳng $x+2y-3=0$ có một véc-tơ pháp tuyến là $\overrightarrow{n}=(1;2)$.
	} 
\end{ex}

% nguồn: MaHoaNH_TracNghiem_Lop10
\begin{ex}%[0H4Y1-1]%[ma:15897b71]
	Đường thẳng $\Delta\colon 2x+3y+2022=0$ có vectơ pháp tuyến  có toạ độ là
	 \choice
		{$\left(2022;2\right)$} 
		{$\left(-3;2\right)$} 
		{$\left(2;2022\right)$} 
		{\True  $\left(2;3\right)$}
	 \loigiai{
		Đường thẳng $\Delta:2x+3y+2022=0$ có vectơ pháp tuyến  là $\left(2;3\right)$.
	 }
\end{ex}

% nguồn: MaHoaNH_TracNghiem_Lop10
\begin{ex}%[0H4B1-1]%[ma:8e279247]
	Trong mặt phẳng $Oxy$, cho đường thẳng $d\colon\heva{&x=-1+3t\\&y=5-2t}$. Véc tơ nào sau đây là  véc-tơ chỉ phương của đường thẳng $d$?
	 \choice
		{$\vec{u}_1=(2;3)$} 
		 {$\vec{u}_2=(3;2)$}
		 {$\vec{u}_3=(-1;5)$}
		 {\True $\vec{u}_4=(3;-2)$}
	\loigiai{ 
		Véc-tơ chỉ phương của  đường thẳng $d$ là $\vec{u}_4=(3;-2)$.
	 }
\end{ex}

% nguồn: MaHoaNH_TracNghiem_Lop10
\begin{ex}%[0H4Y1-1]%[ma:5c7aa6dc]
	Trong mặt phẳng $O x  y$, một vectơ pháp tuyến của đường thẳng $d\colon2 x-3 y+9=0$ là
	\choice 
		{$\vec{n}=(3;-2)$} 
		{$\vec{n}=(-2;-3)$} 
		 {$\vec{n}=(2; 3)$}
		 {\True $\vec{n}=(2;-3)$}
	 \loigiai{
		Vectơ pháp tuyến của  đường thẳng $d\colon2 x-3 y+9=0$ là $\vec{n}=(2;-3)$.
	 }
\end{ex}

% nguồn: MaHoaNH_TracNghiem_Lop10
\begin{ex}%[0H4Y1-1]%[ma:65300251]
	 Trong mặt phẳng $Oxy$, cho đường thẳng $\Delta\colon\heva{&x=5+2 t \\& y=3-t}(t \in \mathbb{R})$. Một véc-tơ chỉ phương của đường thẳng $\Delta$ có tọa độ là
	\choice 
		 {$(2; 1)$}
		 {\True $(2; -1)$}
		{$(1; 2)$} 
		 {$(5; 3)$}
	 \loigiai{
		 Đường thẳng $\Delta\colon\heva{&x=5+2t \\ &y=3-t}(t \in \mathbb{R})$ có một véc-tơ chỉ phương là $\vec{u}=(2;-1)$.
	 }
\end{ex}

% nguồn: MaHoaNH_TracNghiem_Lop10
\begin{ex}%[0H4Y1-1]%[ma:9779bc64]
	Trong mặt phẳng $Oxy$, đường thẳng qua điểm $M(4; 5)$ và nhận $\vec{u}=(2; 3)$ làm  véc tơ chỉ phương có phương trình tham số là
	\choice 
		{$\heva{&x=2+4t \\ &y=3+5t}(t  \in \mathbb{R})$}
		{\True $\heva{&x=4+2t \\  &y=5+3t}(t \in \mathbb{R})$}
		{$\heva{&x=4+5t  \\ &y=2+3t}(t \in \mathbb{R})$}
		{$\heva{&x=2+t \\ &y=5+3t}(t \in \mathbb{R})$} 
	\loigiai{ 
		 Đường thẳng qua điểm $M(4; 5)$ và nhận $\vec{u}=(2; 3)$ làm véc-tơ chỉ phương có phương trình tham số là $\heva{&x=4+2t \\& y=5+3t}(t \in \mathbb{R})$.
	} 
\end{ex}

% nguồn: MaHoaNH_TracNghiem_Lop10
\begin{ex}%[0H4Y1-1]%[ma:52145d53]
	Trong mặt phẳng $Oxy$, cho đường thẳng $\Delta \colon \left\{\begin{aligned}
	 & x=-1+3t\\ & y=4-2t\\
\end{aligned}\right.$. Véc-tơ nào sau đây là véc-tơ chỉ phương của đường thẳng $\Delta$?
 \choice
	 {$\overrightarrow{u}=(3;2)$}
	 {$\overrightarrow{u}=(-1;4)$}
	 {\True $\overrightarrow{u}=(3;-2)$}
	{$\overrightarrow{u}=(-1;-2)$} 
\loigiai{ 
	Đường thẳng $\Delta$ có véc-tơ  chỉ phương là $\overrightarrow{u}=(3;-2)$.
 }
\end{ex}

% nguồn: MaHoaNH_TracNghiem_Lop10
\begin{ex}%[0H4B1-1]%[ma:db2033a3]
	Cho tam  giác $ABC$ có tọa độ ba đỉnh lần lượt là $(2;3)$, $(5;4)$, $(2;2)$. Tọa độ trọng tâm $G$ của tam giác có tọa độ là
	\choice 
		 {$(1;1)$}
		{$(2;2)$} 
		{\True  $(3;3)$}
		 {$(4;4)$}
	 \loigiai{
		 Tọa độ trọng tâm $G$ của tam giác có tọa độ là
		\begin{equation*}
			\heva{&x_G =  \dfrac{x_A+x_B+x_C}{3}=\dfrac{2+5+2}{3}=3\\&y_G = \dfrac{y_A+y_B+y_C}{3}=\dfrac{3+4+2}{3}=3.}
		\end{equation*}
	} 
\end{ex}


%%%%%%%%%%%%%%%% ĐÚNG/SAI %%%%%%%%%%%%%%%%

% nguồn: MaHoaNH_DungSai_Lop10
\begin{ex}%[0H4N1-1]%[ma:4a325b60]
	Với mọi góc $\alpha$ thoả  mãn $0^{\circ} < \alpha < 90^{\circ}$, ta có
	 \choiceTF
		{\True $\sin \alpha  +2\cos \alpha >0$}
		{$\sin \left(90^{\circ}+\alpha\right)  \cdot \cos \alpha <0$}
		{$\tan  \alpha \cdot \cot \alpha <0$}
		{\True $\tan \left(180^{\circ}-\alpha\right) \cdot \cot \left(180^{\circ}-\alpha\right)  >0$}
	\loigiai{ 
		Với mọi góc $\alpha$ thoả mãn $0^{\circ} < \alpha  < 90^{\circ}$, ta có
		\begin{itemchoice}
			\itemch $\sin \alpha >0$, $\cos \alpha >0$ nên $\sin \alpha +2\cos \alpha  >0$.
			 \itemch $\cos \alpha >0$, $\sin \left(90^{\circ}+\alpha\right)>0$ nên $\sin \left(90^{\circ}+\alpha\right) \cdot \cos \alpha >0$.
			\itemch $\tan \alpha >0$, $\cot \alpha >0$ nên $\tan \alpha \cdot \cot \alpha >0$. 
			\itemch $\tan \left(180^{\circ}-\alpha\right)= -\tan \alpha <0$, $ \cot \left(180^{\circ}-\alpha\right)=-\cot \alpha  <0$ nên $\tan \left(180^{\circ}-\alpha\right) \cdot \cot \left(180^{\circ}-\alpha\right) >0$.
		\end{itemchoice}
	 }
\end{ex}

% nguồn: MaHoaNH_DungSai_Lop10
\begin{ex}%[0H4N1-1]%[ma:8a8dfc2b]
	Với mọi góc $\alpha$ thoả mãn $0^{\circ} < \alpha <  90^{\circ}$, ta có
	 \choiceTF
		{\True $\sin \alpha  >0$}
		{$\cos  \alpha <0$}
		{$\tan  \alpha <0$}
		{\True  $\cot \alpha >0$}
	\loigiai{ 
		Với mọi góc $\alpha$ thoả mãn $0^{\circ} < \alpha < 90^{\circ}$, ta  có
		\begin{itemchoice}
			\itemch $\sin \alpha  >0$.
			\itemch $\cos \alpha  >0$.
			\itemch  $\tan \alpha >0$.
			\itemch $\cot  \alpha >0$.
		\end{itemchoice}
	 }
\end{ex}


%%%%%%%%%%%%%%%% TỰ LUẬN %%%%%%%%%%%%%%%%

% nguồn: MaHoaNH_TuLuan_Lop10
\begin{bt}%[0H4Y1-1]%[0H4B1-2]%[ma:ce6d0924]
	Trong mặt phẳng tọa độ $Oxy$,  cho hai điểm $A(2;2)$ và $B(4;-2)$.
	\begin{enumerate}
		\item  Tìm tọa độ véc-tơ $\overrightarrow{AB}$.
		\item Viết phương trình tổng quát  của đường thẳng $AB$.
	\end{enumerate}
	 \loigiai{
	\begin{enumerate}
		\item Ta  có $\overrightarrow{AB}=(2;-4)$.
		\item Đường thẳng  $AB$ nhận véc-tơ $\overrightarrow{AB}=(2;-4)$ làm véc-tơ chỉ phương.\\
		Chọn véc-tơ pháp tuyến của đường thẳng  $AB$ là $\overrightarrow{n}=(2;1)$.\\
		Đường thẳng $AB$  đi qua điểm $A(2;2)$ nên phương trình tổng quát của đường thẳng $AB$ là
		$$2(x-2)+1(y-2)=0  \Leftrightarrow 2x+y-6=0.$$
	\end{enumerate}
	} 
\end{bt}
